\documentclass[11pt,letterpaper]{article}
\usepackage[T1]{fontenc}
\usepackage[utf8]{inputenc}
\usepackage{lmodern,microtype}
\usepackage[margin=0.79in,headheight=15pt]{geometry}
\usepackage{parskip,enumitem,longtable,tabularx,booktabs,array,xcolor,fancyhdr,hyperref}
\definecolor{navy}{HTML}{17314C}
\definecolor{teal}{HTML}{087F78}
\hypersetup{colorlinks=true,linkcolor=navy,urlcolor=teal}
\setlength{\emergencystretch}{3em}
\setlist[itemize]{leftmargin=*,nosep,topsep=5pt}
\setlist[enumerate]{leftmargin=*,itemsep=3pt}
\renewcommand{\arraystretch}{1.22}
\pagestyle{fancy}\fancyhf{}
\fancyhead[L]{\small Enterprise Data Modeling}\fancyhead[R]{\small Documentation Suite}
\fancyfoot[C]{\thepage}
\newcommand{\srcbook}{\textit{The Data Model Resource Book}, revised edition, Volume 2}
\newcommand{\ResponseLine}{\par\vspace{1.05em}\noindent\rule{\linewidth}{0.4pt}\par}
\newcommand{\ResponseArea}{\ResponseLine\ResponseLine\ResponseLine}
\begin{document}
\begin{center}
{\LARGE\bfseries\color{navy} Schema Migration and Rollback Planning}\par
\vspace{0.25em}{\large Guided workbook | conceptual and operational analysis}\par
\vspace{0.3em}{\small Original engineering study material \textbullet{} 2026 edition}
\end{center}
\vspace{0.8em}
\section{Workbook objective}
Complete independently before looking at any model answer. Record assumptions, counterexamples and a proof of each chosen invariant. This is supplementary material; no exercise copies a worked example from the source book.


\section{Problem A: column split}
Given: Existing API consumers read full\_name while new consumers expect given\_name and family\_name.

Tasks: Plan additive fields, backfill, compatibility projection, writer changes, monitoring and contract phase.

\subsection{Your working}
\ResponseArea\ResponseArea
\subsection{Verification notes}
\ResponseArea

\section{Problem B: constraint rollout}
Given: A table has historical null customer\_id rows and the design now requires NOT NULL.

Tasks: Specify data-quality detection, remediation ownership and safe constraint rollout; include rollback decisions.

\subsection{Your working}
\ResponseArea\ResponseArea
\subsection{Verification notes}
\ResponseArea

\section{Problem C: traffic regression}
Given: A new composite index helps one query but hurts write throughput under load.

Tasks: Design representative benchmarks, acceptance budget and operational rollback thresholds.

\subsection{Your working}
\ResponseArea\ResponseArea
\subsection{Verification notes}
\ResponseArea

\section{Self-assessment}
\begin{itemize}
\item Can another reviewer reproduce the decomposition or integrity argument?
\item Did the model document what would falsify its assumptions?
\item Can every proposed invariant be linked to executable or review evidence?
\end{itemize}

\vfill
\hrule
{\footnotesize\textbf{Provenance.} Supplementary engineering practice: original content; the cited book provides domain-modeling context.}
\end{document}
